\bibitem{icd11} World Health Organization, \emph{International Classification of Diseases
11th Revision (ICD-11)} --- Gaming disorder (6C51), 2019.
\bibitem{igds9sf} H. M. Pontes and M. D. Griffiths, ``Measuring DSM-5 Internet Gaming
Disorder: Development and validation of a short psychometric scale,'' \emph{Computers in
Human Behavior}, vol.\ 45, pp.\ 137--143, 2015.
\bibitem{igdslatam2024} J. A. Bland\'on-Casta\~no, D. X. Puerta-Cort\'es, A. Carmona,
J. Ivern, L. W. Vilca, X. Carbonell, and T. Caycho-Rodr\'iguez, ``Cross-cultural invariance of
the Internet Gaming Disorder Scale--Short-Form (IGDS9-SF) in seven Latin American countries,''
\emph{Int.\ J.\ Mental Health and Addiction}, vol.~23, no.~6, pp.~4690--4715, 2024;
open dataset: \texttt{doi:10.17605/OSF.IO/GWCSN}.
\bibitem{shap2017} S. M. Lundberg and S.-I. Lee, ``A unified approach to interpreting
model predictions,'' in \emph{Advances in Neural Information Processing Systems
(NeurIPS)}, 2017.
\bibitem{conda2021} H. Weld \emph{et al.}, ``CONDA: a CONtextual Dual-Annotated dataset
for in-game toxicity understanding and detection,'' in \emph{Findings of ACL-IJCNLP},
2021.
\bibitem{davidson2017} T. Davidson, D. Warmsley, M. Macy, and I. Weber, ``Automated hate
speech detection and the problem of offensive language,'' in \emph{Proc.\ ICWSM}, 2017.
\bibitem{jigsaw2018} Jigsaw / Conversation AI, ``Toxic Comment Classification
Challenge'' (Wikipedia talk-page comments with crowd-sourced toxicity annotations),
Kaggle, 2018.
\bibitem{hasoc2019} T. Mandl, S. Modha, P. Majumder, D. Patel, M. Dave, C. Mandlia,
and A. Patel, ``Overview of the HASOC track at FIRE 2019: Hate speech and offensive
content identification in Indo-European languages,'' in \emph{Proc.\ FIRE}, 2019.
\bibitem{svarah2023} T. Javed, S. Joshi, V. Nagarajan, S. Sundaresan, J. Nawale,
M. Raman, K. Kumar, P. Kumar, and M. M. Khapra, ``Svarah: Evaluating English ASR
systems on Indian accents,'' in \emph{Proc.\ INTERSPEECH}, 2023.
\bibitem{detoxify2020} L. Hanu and Unitary team, ``Detoxify: toxic comment
classification models,'' open-source software,
\texttt{github.com/unitaryai/detoxify}, 2020.
\bibitem{vosk} Alpha Cephei, ``Vosk speech recognition toolkit,'' open-source
software, \texttt{alphacephei.com/vosk}.
\bibitem{bohra2018} A. Bohra, D. Vijay, V. Singh, S. S. Akhtar, and M. Shrivastava,
``A dataset of Hindi-English code-mixed social media text for hate speech
detection,'' in \emph{Proc.\ NAACL PEOPLES Workshop}, 2018.
\bibitem{studentlife2014} R. Wang, F. Chen, Z. Chen, T. Li, G. Harari, S. Tignor,
X. Zhou, D. Ben-Zeev, and A. T. Campbell, ``StudentLife: Assessing mental health,
academic performance and behavioral trends of college students using smartphones,''
in \emph{Proc.\ ACM UbiComp}, 2014.
\bibitem{wikihi} Wikimedia Foundation, ``Hindi Wikipedia,'' 20231101 snapshot
(CC-BY-SA), accessed via the \texttt{wikimedia/wikipedia} dataset, 2026.
\bibitem{ravdess2018} S. R. Livingstone and F. A. Russo, ``The Ryerson Audio-Visual
Database of Emotional Speech and Song (RAVDESS),'' \emph{PLoS ONE}, vol.\ 13, no.\ 5,
2018.
\bibitem{cremad2014} H. Cao, D. G. Cooper, M. K. Keutmann, R. C. Gur, A. Nenkova, and
R. Verma, ``CREMA-D: Crowd-sourced emotional multimodal actors dataset,'' \emph{IEEE
Trans.\ Affective Computing}, vol.~5, no.~4, 2014.
\bibitem{emodb2005} F. Burkhardt, A. Paeschke, M. Rolfes, W. F. Sendlmeier, and B. Weiss,
``A database of German emotional speech,'' in \emph{Proc.\ Interspeech}, 2005.
\bibitem{urdu2018} S. Latif, A. Qayyum, M. Usman, and J. Qadir, ``Cross lingual speech
emotion recognition: Urdu vs.\ Western languages,'' in \emph{Proc.\ Int.\ Conf.\ on
Frontiers of Information Technology (FIT)}, 2018.
\bibitem{mobsf} A. Abraham \emph{et al.}, ``Mobile Security Framework (MobSF),''
open-source software, \texttt{github.com/MobSF/Mobile-Security-Framework-MobSF}.
\bibitem{huang2024} Y. Huang, R. Wu, Y. Huang, Y. Xiang, and W. Zhou,
``Investigating mechanisms of Internet Gaming Disorder and developing intelligent
monitoring models using artificial intelligence technologies,'' \emph{BMC Public
Health}, 2024.
\bibitem{strojny2024} P. Strojny, K. Kapela, N. Lipp, and S. Sikstr\"om, ``Use of four
open-ended text responses to help identify people at risk of gaming disorder,'' \emph{JMIR
Serious Games}, vol.\ 12, e56663, 2024.
\bibitem{ergin2025} D. A. Ergin and C. A. Essau, ``Family factors and Internet Gaming
Disorder among adolescents: A systematic review,'' \emph{Int.\ J.\ Developmental
Science}, 2025.
\bibitem{podsakoff2003} P. M. Podsakoff, S. B. MacKenzie, J.-Y. Lee, and N. P. Podsakoff,
``Common method biases in behavioral research: A critical review of the literature and
recommended remedies,'' \emph{J. Applied Psychology}, vol.~88, no.~5, pp.~879--903, 2003.
\bibitem{kroenke2003} K. Kroenke, R. L. Spitzer, and J. B. W. Williams, ``The Patient Health
Questionnaire-2: Validity of a two-item depression screener,'' \emph{Medical Care},
vol.~41, no.~11, pp.~1284--1292, 2003.
\bibitem{cohen1988} J. Cohen, \emph{Statistical Power Analysis for the Behavioral Sciences},
2nd ed. Hillsdale, NJ: Lawrence Erlbaum, 1988.
